\chapter{1976:William Lipscomb}

\label{chap:chemistry1976}

The Nobel Prize in Chemistry for the year 1976 was awarded to William Lipscomb.

\textbf{Motivation:}

for his studies on the structure of boranes illuminating problems of chemical bonding

\section{William Lipscomb}

\subsection*{Early Life and Education}

William Lipscomb was born on 1919-12-09 in Cleveland, Ohio, USA.

\subsection*{Career and Major Research}

Lipscomb studied at the University of Kentucky and received his doctorate from the California Institute of Technology in 1946. He taught at the University of Minnesota and then at Harvard University, where he was professor of chemistry. He determined the structures of the boranes, including diborane and the higher boron hydrides, and developed theories of their electron-deficient bonding.

\subsection*{Nobel Prize-winning Work}

The work for which William Lipscomb was awarded the Nobel Prize in Chemistry in 1976 was described by the Nobel Committee as: "for his studies on the structure of boranes illuminating problems of chemical bonding".

Lipscomb used X-ray and electron diffraction to determine the structures of the boranes and showed how these electron-deficient molecules are held together by multicenter bonds.

\subsection*{Significance and Impact}

His structural work clarified the nature of the chemical bond in compounds that cannot be described by ordinary two-center two-electron bonds.

\subsection*{Later Career and Honors}

Lipscomb remained at Harvard University until his retirement. He died on 2011-04-14 in Cambridge, Massachusetts, USA.

\subsection*{Legacy}

Lipscomb's borane structures and bonding theories underpin modern boron cluster chemistry.

\subsection*{Modern Developments}

Lipscomb's boron hydride studies revealed three-center, two-electron bonds and founded the modern chemistry of electron-deficient compounds. His methods extended to the elucidation of enzyme mechanisms, including carboxypeptidase. Boron chemistry is now used in drug discovery, including boron neutron capture therapy for cancer, and in the synthesis of complex organic molecules.

\subsection*{Selected Publications}

"Boron Hydrides" (1963)

"The Boranes and Their Relatives" (Nobel Lecture, 1976)

\clearpage

\section*{Chapter Summary}

The 1976 prize honored William Lipscomb for his structural studies of the boranes. His work illuminated the nature of chemical bonding in electron-deficient molecules.

